\section*{ПОСТАНОВКА ЗАДАЧІ}

\textbf{Дано:} цифрові растрові зображення різної роздільної здатності та базова теоретична концепція глобальної поліноміальної інтерполяції зображень на нерівномірних сітках взята за основу~\cite{occorsio2022lagrange}.

\textbf{Завдання:}
\begin{enumerate}
	\item Ознайомитися з існуючими методами просторового масштабування зображень.
	\item Адаптувати математичну модель цифрового зображення на нерівномірних сітках для випадку вузлів Чебишева другого роду.
	\item Поширити другу інтерполяційну барицентричну формулу по вузлах Чебишева на двовимірний випадок.
	\item Розробити програмне забезпечення для наочної демонстрації запропонованого методу.
	\item Провести автоматизоване тестування та об'єктивний порівняльний аналіз результатів роботи розробленого алгоритму з класичними методами масштабування зображень.
	\item Результати досліджень оформити у вигляді дипломної роботи згідно з встановленими вимогами.
\end{enumerate}

Постановку задачі було сформульовано за участі наукового консультанта доктора фізико-математичних наук Василика Віталія Богдановича.
\newpage